\documentclass[12pt, a4paper]{article}
\usepackage[utf8]{inputenc}
\usepackage[spanish]{babel}
\usepackage{amsmath, amssymb}
\usepackage{geometry}
\usepackage{enumitem}
\geometry{top=2.5cm, bottom=2.5cm, left=2.5cm, right=2.5cm}

\title{Resolución de Ejercicio: Técnicas de Conteo}
\author{}
\date{}

\begin{document}

\maketitle

\section*{Enunciado}
17. En una bolsa hay cuatro bolas numeradas con los dígitos 2, 4, 7 y 9. Elegimos una bola de la bolsa y anotamos su número. La bola extraída se introduce en la bolsa. Se elige una segunda bola y se anota su número. La bola extraída se vuelve a introducir en la bolsa. Finalmente se elige una tercera bola y se anota su número. ¿Cuántos números de tres cifras podemos obtener? Ejemplo: se puede obtener el número 222.

\vspace{0.5cm}
\hrule
\vspace{0.5cm}

\section*{Solución Paso a Paso}

\subsection*{1. Análisis e identificación del modelo}
Para determinar de cuántas formas se pueden formar estos números de tres cifras, analizamos las características del experimento estadístico:
\begin{itemize}[label=\checkmark]
    \item \textbf{¿Importa el orden?} Sí. Al construir un número de varias cifras, el orden de los dígitos altera el resultado final (por ejemplo, el número 247 es distinto al 427)[cite: 1].
    \item \textbf{¿Intervienen todos los elementos?} No. La bolsa contiene un total de $m=4$ opciones (los dígitos 2, 4, 7, 9), de los cuales solo seleccionaremos $n=3$ para formar cada número[cite: 1].
    \item \textbf{¿Se pueden repetir los elementos?} Sí. El enunciado indica explícitamente que la bola extraída siempre ``se vuelve a introducir en la bolsa'' antes de la siguiente extracción. Esto implica que hay reposición y un mismo dígito puede salir más de una vez (como en el ejemplo 222)[cite: 1].
\end{itemize}

Como el orden importa, no se usan todos los elementos a la vez, y existe reposición (repetición permitida), el modelo combinatorio compuesto que corresponde aplicar es el de \textbf{Variaciones con repetición}[cite: 1].

\subsection*{2. Planteo matemático y cálculo}
La fórmula estadística para calcular las variaciones con repetición de $n$ elementos elegidos de un total de $m$ opciones es[cite: 1]:
\[ V'_{(m,n)} = m^n \]

Definimos nuestra cantidad de opciones disponibles $m = 4$ (los cuatro dígitos) y la cantidad de extracciones o posiciones a llenar $n = 3$ (las tres cifras del número), y procedemos a calcular:
\begin{align*}
V'_{(4,3)} &= 4^3 \\
V'_{(4,3)} &= 4 \times 4 \times 4 \\
V'_{(4,3)} &= 64
\end{align*}

Llegaríamos al mismo resultado utilizando la regla de la multiplicación: tenemos 4 opciones posibles para la primera cifra, 4 opciones para la segunda (porque la bola regresó a la bolsa) y 4 opciones para la tercera cifra ($4 \times 4 \times 4 = 64$)[cite: 1].

\subsection*{3. Conclusión}
Al realizar tres extracciones sucesivas con reposición a partir de cuatro dígitos diferentes, se generan 64 arreglos ordenados posibles.

\vspace{0.5cm}

\textbf{Respuesta:} Podemos obtener \textbf{64 números} de tres cifras distintos.

\end{document}
